% from ega4-11-fr.tex
\oldpage[IV-3]{149}

$V=U\cap T'_1$ est ouvert non vide dans $T'_1$. L'anneau $A$, étant artinien,
est un anneau de Jacobson, donc (10.4.6) $Z'_1$ et $T'_1$ sont des préschémas de
Jacobson~; il existe par suite dans $V$ un point $t'_0$ fermé dans $V$, et son
image $z'_0$ dans $Z'_1$ est un point fermé de $Z'_1$ (10.4.7). Soient $f_1$ la
restriction de $f$ à $X_1$ et $f'_1:T'_1\to Z'_1$ la projection canonique~;
$V\cap f_1'^{-1}(z'_0)$ est un ouvert non vide dans
$f_1^{-1}(y)\otimes_{k(y)}k(z'_0)$, et comme ce dernier préschéma est plat sur
$f_1^{-1}(y)$, un point maximal $t'_1$ de $V\cap f_1'^{-1}(z'_0)$ est
nécessairement au-dessus de l'unique point maximal $x$ de $X_1$ (2.3.4).
Enfin, $k(z'_0)$ est une extension finie de $k$ (\textbf I, 6.4.2) et
l'homomorphisme $A=\mathcal O_{Y,y}\to A'=\mathcal O_{Z',z'}$ se factorise en
$\mathcal O_{Y,y}\to\mathcal O_{Z',z'_0}\to\mathcal O_{Z',z'}$, donc son noyau
est contenu dans celui de $\mathcal O_{Y,y}\to\mathcal O_{Z',z'_0}$. Ceci
achève la réduction annoncée.

\medskip
\noindent\textbf{II)} \emph{Réduction au cas où les $A'_\alpha$ sont en nombre
fini et sont des $A$-algèbres finies.} --- Soit $\mathfrak m'_\alpha$ l'idéal
maximal de $A'_\alpha$~; comme $A'_\alpha$ est un anneau local noethérien,
l'intersection des $\mathfrak m'^n_\alpha$ est $0$ (\textbf I, 7.3.5)~;
l'intersection des $u_\alpha^{-1}(\mathfrak m'^n_\alpha)$, pour tous les indices
$n$ et $\alpha$, est donc égale à l'intersection des noyaux des $u_\alpha$, donc
est réduite à $0$ par l'hypothèse (i). Puisque $A$ est artinien, il y a un nombre
fini de ces idéaux dont l'intersection est déjà $0$~; notons-les
$u_i^{-1}(\mathfrak m_i'^{\,n_i})$ $(1\leq i\leq r)$. Comme les $A'_i$ sont
noethériens, les $\mathfrak m_i'^{\,h}/\mathfrak m_i'^{\,h+1}$ sont des
$(A'_i/\mathfrak m'_i)$-modules de longueur finie, et comme
$A'_i/\mathfrak m'_i$ est un $k$-espace vectoriel de rang fini, on voit que
$A''_i=A'_i/\mathfrak m_i'^{\,n_i}$ est une $A$-algèbre finie et un anneau local
artinien. La réduction annoncée résulte donc de (2.1.4).

\medskip
\noindent\textbf{III)} \emph{Fin de la démonstration.} --- Supposons désormais
que les $A'_i$ $(1\leq i\leq r)$ soient en nombre fini et soient des
$A$-algèbres finies. Pour tout $i$, le corps résiduel $k_i$ de $A'_i$ est une
extension finie de $k$~; utilisant l'hypothèse (ii), on en conclut que l'image
réciproque de $x$ dans $X'_i$ est réduite au seul point $x'_i$ (4.3.2). Soient
$Y'$ le préschéma somme des $Y'_i$, $X'=X\times_Y Y'$, somme des $X'_i$,
$\mathcal F'=\mathcal F\otimes_Y Y'$. L'hypothèse entraîne que
$\mathcal F'$ est $Y'$-plat aux points de l'image réciproque de $x$ par la
projection $p:X'\to X$. Comme $p$ est de type fini, il existe par suite un
ouvert $U'\supset p^{-1}(x)$ tel que $\mathcal F'$ soit $Y'$-plat aux points de
$U'$ (11.1.1). En outre, le morphisme $Y'\to Y$ est fini puisque les $A'_i$
sont des $A$-algèbres finies~; donc $p$ est un morphisme fini (\textbf{II},
6.1.5), par suite fermé (\textbf{II}, 6.1.10), et il existe donc un voisinage
affine $U$ de $x$ dans $X$ tel que $p^{-1}(U)\subset U'$. Soient $B$ et $A'$
les anneaux des schémas $U$ et $Y'$ ($A'$ étant le composé direct des $A'_i$)~;
remplaçant $X$ par $U$, on a donc $\mathcal F=\widetilde M$, où $M$ est un
$B$-module et, par hypothèse, $M'=M\otimes_A A'$ est un $A'$-module plat
(2.1.2)~; comme l'homomorphisme $A\to A'$ est injectif par construction, on
peut appliquer (11.4.3), qui prouve que $M$ est un $A$-module plat, C.Q.F.D.
\end{proof}

\begin{proposition}[11.4.12]
\label{IV.11.4.12-fr}
Soient $A$ un anneau local artinien de corps résiduel $k$,
$Y=\operatorname{Spec}(A)$, $f:X\to Y$ un morphisme localement de type fini,
$\mathcal F$ un $\mathcal O_X$-Module cohérent. Soient $(A'_\alpha)$ une famille
d'anneaux locaux, et pour chaque $\alpha$, $u_\alpha:A\to A'_\alpha$ un
homomorphisme d'anneaux. Posons $Y'_\alpha=\operatorname{Spec}(A'_\alpha)$,
$X'_\alpha=X\times_Y Y'_\alpha$, $f'_\alpha=f_{(Y'_\alpha)}$,
$\mathcal F'_\alpha=\mathcal F\otimes_A A'_\alpha$. Soit $x$ un point de $X$,
et supposons vérifiées les conditions suivantes~:
\begin{enumerate}[label=(\roman*)]
  \item L'intersection des noyaux des $u_\alpha$ est réduite à $0$.
  \item Pour tout $\alpha$, $\mathcal F'_\alpha$ est $f'_\alpha$-plat en tous
  les points $x'_\alpha\in X'_\alpha$ dont les projections respectives dans
  $X$ et $Y'_\alpha$ sont $x$ et le point fermé $y'_\alpha$ de $Y'_\alpha$.
\end{enumerate}
Alors $\mathcal F$ est $f$-plat au point $x$.
\end{proposition}

\begin{proof}
% IV3_B22_P149_ANCHOR

